\chapter{Summary}
\label{cha:Summary}

Potato (\textit{Solanum tuberosum L.}) currently stands as the world's third most important food crop after rice and wheat. Potato provides essential nutrition and calories to billions of people globally. Its tubers are prized not only for their versatility in culinary applications but also for their nutritional value, offering a rich source of carbohydrates, vitamins, minerals, and antioxidants. Beyond its role as a staple food, potato serves as a critical raw material for industrial applications including starch production, processed foods, and bio-based materials. Despite its essential role in international food security, genetic improvement of potato faces formidable challenges and is constrained by multiple biological and market-related factors. The crop's tetraploid nature, clonal propagation system, and complex genetic architecture complicate breeding efforts and negatively affect innovation. Furthermore, potato's susceptibility to numerous pests and diseases, coupled with its sensitivity to environmental stresses, necessitates continuous genetic innovation. These realities coupled with long breeding cycles characteristic of clonally propagated crops hamper the creation of novel potato varieties. Hybrid and diploid breeding systems have been proposed as a potential solution which could circumvent some of the issues facing traditional clonal breeding. Hybrid breeding systems are inherently true seed based which allow for faster propagation and multiplication rates in a breeding pipeline. True potato seed can be leveraged in replicated early yield and multi-environmental trialling resulting in higher quality evaluations of important traits such as yield. This in turn increases the accuracy and potentially the speed of selection earlier in the breeding cycle. This further enables the utilization of methodologies such as quantitative trait locus (QTL) mapping and genomic prediction (GP) which depend on high quality phenotypic data. This transforms the traditional potato breeding into an inherently data-driven pipeline. However, this is all contingent on several genetic assumptions in how hybrid potato behaves. This includes a sufficient level of heterotic vigour, high heritabilities among complex traits, and consistent uniformity in the F1 hybrids. Characterising these factors are crucial for any implementation of hybrid potato. This thesis assesses these questions in the following form:

In \textbf{Chapter 1} I begin with a brief introduction of the evolutionary and recent cultivated history of potato. This is followed by a justification for genetic improvement in potato and where I lay out current popular breeding methods for clonal potato. This includes an examination of specific factors of potato's biology which directly inform breeding. Specifically, how it is hampered by potato's tetraploidy, obligate clonal propagation, and deleterious mutation load. I then focus on hybrid breeding methodology and assess where it can potentially augment genetic gain using diploid potato. However, to properly evaluate whether hybrid breeding can do what it has promised, a proper quantitative evaluation must be made using extant material. This chapter concludes by describing this knowledge gap and the aims of this thesis.

A central component of hybrid breeding is the ability to conduct multi-environmental trials (MET) and be able to retrieve accurate estimates of merit (e.g. best linear unbiased predictor) and high enough heritabilities to warrant use in a breeding program. \textbf{Chapter 2} examines this core premise for the first time in hybrid potato utilizing METs from an active breeding program. Me and my coauthors begin with a thorough description of the trial data as it is one of the largest hybrid potato METs published. We showcased a two-step statistical evaluation of the MET where intra-trial effects were first estimated and fitted for all observations. This was followed by a secondary multi-trial model where the spatially-resolved yield data was used to estimate multiple genetical parameters, namely general and specific combining ability (GCA \& SCA respectively). We found that the modelling yield on the GCA of the parents was sufficient for all yield components studied. All traits had significant SCA effects as well indicative of non-additive genetic effects not captured by the parental effects alone. We also examined the genetic architecture of these yield components via a multivariate linear mixed model methodology. This procedure allowed us to delineate these genetic correlations into additive and non-additive components. These results identified average tuber volume as a potentially useful trait for long-term yield gains in a hybrid potato program.

Because \textbf{Chapter 2} showed that hybrid potato \textit{behaves} like a hybrid with respect to genetic parameters, it was natural to extend the use of this trial data. Hybrid breeding is often improved through the use predictive methodologies such as genomic prediction which can increase selection accuracy, offer pre-selection capabilities, and lessen breeding cycle length. Because of this, examining the current capacity of genomic prediction in this hybrid potato dataset is a crucial exercise. To do this, \textbf{Chapter 3} built on the two-step phenotypical modelling approach of \textbf{Chapter 2} and built on it through the application of genetic marker information. To assess the quality of the molecular marker information, we conducted a thorough marker filtering analysis along with an examination of population structure among the parental lines. We found moderate population structure among the parental lines indicative of minor parental grouping, but no proper pool structure. We then created two genomic models mirroring the structure of the  From here, I fit several genomic models and examine their genetic parameters as well as the model performance for each yield component of interest. Here, I examine the benefits of predicting hybrid yield performance with multiple genetic components over a simpler model with one genetic component.

After having surveyed a few genomic models in \textbf{Chapter 3}, I wanted to consider genome-wide prediction models in hybrid potato relative to other classic technologies such as quantitative trait locus (QTL) mapping. \textbf{Chapter 4} examines this question in the context of several selection scenarios where I compare the selection efficiency of a genomic selection scenario relative to marker-assisted selection (MAS). I consider the efficiency of these scenarios both in terms of genetic gain as well as based upon unit gain per cost.

In my final experimental chapter, I examined the utility of multiallelic molecular marker information for genome-wide predictive models. \textbf{Chapter 5} examines multiple forms of multiallelic marker information, identity-by-state (IBS) marker haplotags and estimated founder origin identity-by-descent (IBD) variants. This chapter examines the benefit of either of these methods relative to traditional biallelic single nucleotide polymorphisms (SNPs) in the context of their effects of model performance when used as a predictor.

I conclude this thesis with a general discussion and retrospective in \textbf{Chapter 6}. I take the time to examine the implications of these results for future hybrid potato breeding and locate questions which are still unaddressed at this point in time.
